% ---------------------------------------------------------------------------
% The Level Decomposition for Trajectory-Balance Samplers
% Compiles with any standard TeX Live / Overleaf installation:
%     pdflatex tb_level_decomposition.tex   (twice, for cross-references)
% ---------------------------------------------------------------------------
\documentclass[11pt]{article}

\usepackage[a4paper,margin=1.15in]{geometry}
\usepackage{amsmath,amssymb,amsthm}
\usepackage{microtype}
\usepackage{booktabs}
\usepackage{xcolor}
\usepackage[colorlinks=true,linkcolor=black,citecolor=black,urlcolor=blue]{hyperref}

% --- theorem environments --------------------------------------------------
\theoremstyle{plain}
\newtheorem{theorem}{Theorem}
\newtheorem{proposition}[theorem]{Proposition}
\newtheorem{lemma}[theorem]{Lemma}
\newtheorem{corollary}[theorem]{Corollary}
\theoremstyle{definition}
\newtheorem{definition}[theorem]{Definition}
\theoremstyle{remark}
\newtheorem{remark}[theorem]{Remark}

% --- notation --------------------------------------------------------------
\newcommand{\KL}[2]{D_{\mathrm{KL}}\!\left(#1 \,\middle\|\, #2\right)}
\newcommand{\E}[2]{\mathbb{E}_{#1}\!\left[#2\right]}
\newcommand{\PF}{P_{F}}
\newcommand{\PB}{P_{B}}
\newcommand{\Zstar}{Z^{\star}}
\newcommand{\JF}{J_{F}}
\newcommand{\JB}{J_{B}}
\newcommand{\ent}{\mathbb{H}}

\title{\textbf{The Level Decomposition for Trajectory-Balance Samplers}\\[0.4em]
\large What the forward and backward log-partition estimates actually measure}
\author{}
\date{}

\begin{document}
\maketitle
\vspace{-3.5em}

\begin{abstract}
\noindent
When training a GFlowNet with the trajectory-balance (TB) objective one routinely
monitors two quantities that both look like estimates of $\log Z$: one computed
from on-policy forward rollouts, one from backward rollouts out of a replay or
prior buffer. They disagree, and the size of the disagreement---call it
$\Delta$---is often used as a convergence criterion. This note shows that
$\Delta$ is \emph{exactly} a signed combination of three Kullback--Leibler
divergences,
\[
  \Delta \;=\; A + B - C,
\]
measuring respectively the policy's mode-seeking error, its mass-covering error,
and the quality of the buffer. Four practical consequences follow, the sharpest
being that $\Delta$ is \emph{not} a valid optimisation target: it can be driven
toward zero by making the buffer worse. We also derive the one-sided constraint
that bounds how low the learned normaliser may be placed before backward
training is asked for something impossible. Everything is elementary---the only
tools are Jensen's inequality and the chain rule for KL divergence---but the
bookkeeping is easy to get wrong, which is the reason for writing it down.
\end{abstract}

\section{Motivation}

A TB sampler carries a learned scalar $\log Z$ (or, in the conditional case, a
learned function $\log Z(c)$) that is supposed to converge to the true log
partition function. In practice one estimates it in two different ways, and the
two estimates disagree persistently. It is tempting to read the disagreement as
``$Z$ has not converged yet'' and to train until it closes.

That reading is wrong, and the purpose of this note is to say precisely what the
disagreement \emph{is}. The short answer: it is a statement about the
\emph{policy} and the \emph{buffer}, and contains no information about $Z$ at
all. $Z$ cannot reduce it, and one of the three terms can be manipulated in a
direction that makes the diagnostic look better while the sampler gets worse.

\section{Setup and notation}

\paragraph{Objects.}
Let $\tau = (s_0 \to s_1 \to \cdots \to s_T)$ denote a complete trajectory from a
fixed initial state $s_0$, and write $x = x(\tau) := s_T$ for its terminal state.
We are given a strictly positive reward $R(x) > 0$ and wish to sample from the
associated Boltzmann target
\[
  \pi(x) \;:=\; \frac{R(x)}{\Zstar},
  \qquad
  \Zstar \;:=\; \sum_{x} R(x).
\]
The model consists of
\begin{itemize}
  \item a \emph{forward policy} $\PF(\tau)$, a normalised distribution over
        complete trajectories, $\sum_\tau \PF(\tau) = 1$;
  \item a \emph{backward policy} $\PB(\tau \mid x)$, normalised over the
        trajectories terminating at $x$, i.e.\ $\sum_{\tau \to x} \PB(\tau\mid x) = 1$
        for every $x$;
  \item a learned scalar $Z$.
\end{itemize}
We write $\PF^{\top}(x) := \sum_{\tau \to x} \PF(\tau)$ for the \emph{terminal
marginal} of the forward policy---the distribution we ultimately care about.

\paragraph{A notational convention that does all the work.}
Given any distribution $q$ over terminal states, let $q\PB$ denote the
distribution over \emph{trajectories} obtained by first drawing $x \sim q$ and
then $\tau \sim \PB(\cdot \mid x)$:
\[
  (q\PB)(\tau) \;:=\; q\bigl(x(\tau)\bigr)\,\PB\bigl(\tau \mid x(\tau)\bigr).
\]
This is normalised, because $\sum_\tau q(x)\PB(\tau\mid x) = \sum_x q(x) \sum_{\tau\to x}\PB(\tau\mid x) = \sum_x q(x) = 1$.
In particular $\pi\PB$ is a perfectly good distribution over trajectories, and it
is the object the forward policy is really trying to match.

\paragraph{The trajectory-balance condition.}
TB seeks parameters such that
\begin{equation}\label{eq:tb}
  Z\,\PF(\tau) \;=\; R\bigl(x(\tau)\bigr)\,\PB\bigl(\tau \mid x(\tau)\bigr)
  \qquad \text{for every } \tau .
\end{equation}
Summing \eqref{eq:tb} over the trajectories terminating at a fixed $x$ gives
$Z\,\PF^{\top}(x) = R(x)$; summing that over $x$ forces $Z = \Zstar$ and hence
$\PF^{\top} = \pi$. So \eqref{eq:tb} is exactly the statement ``the sampler is
correct and $Z$ is the true partition function''.

\paragraph{The importance weight.}
Define, for a trajectory $\tau$,
\begin{equation}\label{eq:logw}
  \log w(\tau) \;:=\; \log R\bigl(x(\tau)\bigr) \;+\; \log \PB\bigl(\tau\mid x(\tau)\bigr) \;-\; \log \PF(\tau).
\end{equation}
Comparing with \eqref{eq:tb}: TB holds exactly if and only if $\log w(\tau) = \log Z$
for \emph{every} $\tau$. Thus $\log w$ is constant at the optimum, and both its
\emph{level} and its \emph{spread} are diagnostics of TB violation. This note is
about the level.

\begin{remark}[Conditional case]
Everything below holds verbatim per condition $c$: replace $R, \pi, \Zstar, Z$ by
$R(\cdot\mid c), \pi(\cdot\mid c), \Zstar(c), Z(c)$ and let all policies be
conditioned on $c$. Nothing in the algebra couples different conditions. In an
implementation the quantities are therefore computed per condition and then
reduced across conditions (by a mean, or by a worst-case quantile).
\end{remark}

\subsection{Two facts we will use twice}

\begin{lemma}[Importance-sampling identity]\label{lem:is}
$\E{\PF}{w} = \Zstar$.
\end{lemma}
\begin{proof}
$\displaystyle \E{\PF}{w} = \sum_\tau \PF(\tau)\,\frac{R(x)\PB(\tau\mid x)}{\PF(\tau)}
 = \sum_x R(x)\!\!\sum_{\tau \to x}\!\!\PB(\tau\mid x) = \sum_x R(x) = \Zstar.$
\end{proof}

\noindent
Note this holds for the \emph{forward} sampling distribution specifically; it is
not true for an arbitrary sampler, and that asymmetry is the source of everything
that follows.

\begin{lemma}[KL chain rule with a shared kernel]\label{lem:chain}
For distributions $q, p$ over terminal states, $\KL{q\PB}{p\PB} = \KL{q}{p}$.
\end{lemma}
\begin{proof}
$\displaystyle \KL{q\PB}{p\PB} = \sum_\tau q(x)\PB(\tau\mid x)\log\frac{q(x)\PB(\tau\mid x)}{p(x)\PB(\tau\mid x)}
= \sum_x q(x) \log\frac{q(x)}{p(x)} = \KL{q}{p},$
since the $\PB$ factors cancel inside the logarithm and $\PB(\cdot\mid x)$ sums to one.
\end{proof}

\section{Two sampling distributions, two levels}

In training we evaluate $\log w$ under two different distributions over
trajectories.

\begin{definition}[The two levels]
Let $\mu$ be the distribution over terminal states held in the buffer. Define
\[
  \JF \;:=\; \E{\PF}{\log w}
  \qquad\text{(forward level)},
  \qquad\qquad
  \JB \;:=\; \E{\mu\PB}{\log w}
  \qquad\text{(backward level)} .
\]
$\JF$ is the empirical mean of $\log w$ over on-policy rollouts; $\JB$ is the
same average taken over buffer draws with backward rollouts.
\end{definition}

\begin{definition}[The three divergences]\label{def:abc}
\[
  A \;:=\; \KL{\PF}{\pi\PB},
  \qquad
  B \;:=\; \KL{\mu\PB}{\PF},
  \qquad
  C \;:=\; \KL{\mu}{\pi}.
\]
All three are non-negative. Informally: $A$ is how far the policy is from the
target (\emph{mode-seeking} error, model in the first argument); $B$ is how badly
the policy fails to cover what the buffer contains (\emph{mass-covering} error,
model in the second argument); $C$ is the quality of the buffer itself.
\end{definition}

\begin{lemma}[Forward level]\label{lem:JF}
$\JF = \log \Zstar - A$. In particular $\JF \le \log\Zstar$ always.
\end{lemma}
\begin{proof}
Substituting $R = \Zstar\pi$ into \eqref{eq:logw},
\[
  \JF = \E{\PF}{\log \frac{\Zstar\,\pi(x)\PB(\tau\mid x)}{\PF(\tau)}}
      = \log\Zstar + \E{\PF}{\log\frac{(\pi\PB)(\tau)}{\PF(\tau)}}
      = \log\Zstar - \KL{\PF}{\pi\PB}. \qedhere
\]
\end{proof}

\begin{remark}
Lemma~\ref{lem:JF} is the usual ELBO/Jensen statement in disguise: by
Lemma~\ref{lem:is} and Jensen, $\JF = \E{\PF}{\log w} \le \log \E{\PF}{w} = \log\Zstar$,
and Lemma~\ref{lem:JF} identifies the gap exactly as $A$.
\end{remark}

\begin{lemma}[Backward level]\label{lem:JB}
$\JB = \log\Zstar - C + B$.
\end{lemma}
\begin{proof}
Write $Q := \mu\PB$. As before,
$\JB = \log\Zstar + \E{Q}{\log\frac{(\pi\PB)(\tau)}{\PF(\tau)}}$.
Insert $Q$ inside the logarithm:
\[
  \E{Q}{\log\frac{\pi\PB}{\PF}}
  = \underbrace{\E{Q}{\log\frac{\pi\PB}{Q}}}_{=\,-\KL{Q}{\pi\PB}}
  \;+\; \underbrace{\E{Q}{\log\frac{Q}{\PF}}}_{=\,\KL{Q}{\PF}\,=\,B}.
\]
By Lemma~\ref{lem:chain}, $\KL{Q}{\pi\PB} = \KL{\mu\PB}{\pi\PB} = \KL{\mu}{\pi} = C$.
\end{proof}

\begin{remark}[$\JB$ is \emph{not} a lower bound on $\log\Zstar$]
Unlike $\JF$, the backward level can exceed $\log\Zstar$: by Lemma~\ref{lem:JB}
it does so whenever $B > C$. This is not a paradox---Lemma~\ref{lem:is} only
guarantees unbiasedness of $w$ under $\PF$. It is a common source of confusion
when the two levels are treated as interchangeable ``estimates of $\log Z$''.
\end{remark}

\section{The master identity}

\begin{theorem}[Level decomposition]\label{thm:master}
With $A,B,C$ as in Definition~\ref{def:abc},
\[
  \boxed{\;\Delta \;:=\; \JB - \JF \;=\; A + B - C\;}
\]
\end{theorem}
\begin{proof}
Subtract Lemma~\ref{lem:JF} from Lemma~\ref{lem:JB}; the $\log\Zstar$ cancels.
\end{proof}

Two sanity checks. If the sampler is perfect ($\PF = \pi\PB$, so $A = 0$) and the
buffer is thermal ($\mu = \pi$, so $C = 0$), then $B = \KL{\pi\PB}{\pi\PB} = 0$ and
$\Delta = 0$. If instead the policy is perfect but the buffer is arbitrary, then
$A=0$ and $B = \KL{\mu\PB}{\pi\PB} = C$, so $\Delta = 0$ again: a bad buffer alone
does not open a gap. Conversely, if the buffer is thermal but the policy is not,
\[
  \Delta \;=\; \KL{\pi\PB}{\PF} + \KL{\PF}{\pi\PB},
\]
the \emph{symmetrised} KL (Jeffreys divergence) between the policy and its
target. So in the well-prepared-buffer regime, $\Delta$ is a genuine
distributional distance, vanishing only when $\PF = \pi\PB$ exactly.

\section{The one-sided constraint on \texorpdfstring{$\log Z$}{log Z}}

The learned normaliser enters through the TB residual
\begin{equation}\label{eq:resid}
  r(\tau) \;:=\; \log Z + \log\PF(\tau) - \log\PB(\tau\mid x) - \log R(x)
  \;=\; \log Z - \log w(\tau),
\end{equation}
whose square is the TB loss. Averaging \eqref{eq:resid} under the backward stream
gives $\E{\mu\PB}{r} = \log Z - \JB$. So the backward residual is centred exactly
when $\log Z = \JB$. How low can that be?

\begin{proposition}[ELBO floor]\label{prop:floor}
Let $\mathrm{ELBO}(\mu) := \E{\mu}{\log R} + \ent(\mu) = \log\Zstar - C$, where
$\ent$ denotes Shannon entropy. Then
\[
  \JB \;\ge\; \mathrm{ELBO}(\mu),
\]
with equality if and only if $\PF = \mu\PB$.
\end{proposition}
\begin{proof}
Immediate from Lemma~\ref{lem:JB}: $\JB - \mathrm{ELBO}(\mu) = B \ge 0$, and
$B = \KL{\mu\PB}{\PF} = 0$ iff $\PF = \mu\PB$. The identity
$\log\Zstar - C = \E{\mu}{\log R} + \ent(\mu)$ is a one-line expansion of
$\KL{\mu}{\pi}$.
\end{proof}

\begin{corollary}[Backward training can be asked for the impossible]\label{cor:impossible}
Suppose $\log Z$ is held fixed (detached) at a value $z_0$. Then the mean backward
residual $\E{\mu\PB}{r} = z_0 - \JB$ can be driven to zero by training $\PF$ only
if $z_0 \ge \mathrm{ELBO}(\mu)$. If $z_0 < \mathrm{ELBO}(\mu)$, no forward policy
whatsoever makes the backward residual mean-zero.
\end{corollary}
\begin{proof}
$\JB$ ranges over $[\mathrm{ELBO}(\mu), \infty)$ as $\PF$ varies: the lower end is
Proposition~\ref{prop:floor}, and $\JB$ can be made arbitrarily large by letting
$\PF$ assign vanishing mass to the buffer's trajectories, since
$B = \E{\mu\PB}{\log(\mu\PB/\PF)} \to \infty$. Setting $z_0 = \JB$ is therefore
achievable exactly when $z_0 \ge \mathrm{ELBO}(\mu)$.
\end{proof}

This is a genuinely \emph{one-sided} constraint, and it is worth stating the
asymmetry plainly: placing $\log Z$ \emph{too high} is harmless---the policy
simply raises $\JB$ to meet it---whereas placing it \emph{below} $\mathrm{ELBO}(\mu)$
demands more forward probability mass on the buffer than normalisation permits.
Under a squared-error loss the optimiser does not stop when asked for the
impossible; it distorts $\PF$ as far as it can in the demanded direction, which
degrades the sampler rather than merely stalling it.

\section{Consequences}

\subsection{$Z$ is a nuisance parameter, not an objective}

$\log Z$ is a single number, but the two streams place it in two different spots,
separated by $\Delta$:
\[
  \E{\PF}{r} = \log Z - \JF,
  \qquad
  \E{\mu\PB}{r} = \log Z - \JB
  = \bigl(\log Z - \JF\bigr) - \Delta .
\]
Choosing $\log Z = \JF$ zeroes the forward residual and leaves the backward
residual at $-\Delta$; choosing $\log Z = \JB$ does the reverse. There is no
choice that zeroes both unless $\Delta = 0$, and by Theorem~\ref{thm:master}
$\Delta$ depends only on $(\PF,\PB,\mu,\pi)$---not on $Z$.

\begin{center}
\emph{The placement of $Z$ decides which stream absorbs the error.\\
It does not decide whether error exists.}
\end{center}

A practical corollary: a controller rule that watches the residual of the branch
that \emph{trains} $Z$ is measuring a tautology. If the forward branch updates
$Z$ and only $Z$, then $\log Z \to \JF$ and the forward residual self-zeroes
regardless of how bad the policy is. The informative quantity is the residual on
the branches that train the \emph{policy}.

\subsection{$\Delta$ is not a valid optimisation target}

By Theorem~\ref{thm:master}, $\Delta = A + B - C$ is \emph{decreasing} in $C$.
Since $C = \KL{\mu}{\pi}$ measures how bad the buffer is, degrading the buffer
reduces $\Delta$. A training loop that churns its own buffer and gates on
$\Delta$ therefore has a direct incentive to make the buffer worse. Whatever
$\Delta$ is used for, it should not be a success criterion; it is a
\emph{handoff diagnostic}---an indicator of whether the placement of $Z$ is
currently contentious.

\subsection{$A$ is the objective}

Of the three terms only $A$ is an unambiguous statement about the thing we want.
Moreover it controls the error in the terminal marginal, which is what is
actually sampled:

\begin{proposition}\label{prop:dpi}
$\KL{\PF^{\top}}{\pi} \;\le\; A$.
\end{proposition}
\begin{proof}
Factor each trajectory distribution into its terminal marginal and the
conditional over trajectories given the terminal state. The chain rule for KL
gives
\[
  A = \KL{\PF}{\pi\PB} = \KL{\PF^{\top}}{\pi} \;+\; \E{x\sim \PF^{\top}}{\KL{\PF(\cdot\mid x)}{\PB(\cdot\mid x)}},
\]
and the second term is non-negative.
\end{proof}

So driving $A$ down drives the sampling error down, with the slack being exactly
the mismatch between the forward posterior over trajectories and the backward
policy. Neither $B$ nor $\Delta$ admits such a bound.

\subsection{$A$ and $B$ require different data}

$A = \KL{\PF}{\pi\PB}$ has the model in the first argument: it is the
\emph{reverse} (mode-seeking) KL, and its gradient is naturally estimated from
samples drawn from $\PF$---that is, from forward rollouts, or from a replay buffer
of previously generated forward trajectories. $B = \KL{\mu\PB}{\PF}$ has the model
in the second argument: it is a \emph{forward} (mass-covering) KL, estimated from
samples drawn from $\mu$---that is, from backward rollouts out of the buffer.

Consequently:
\begin{itemize}
  \item Training on buffer data alone reduces $B$ and leaves $A$ unattended: the
        sampler learns to cover what it is shown and never sharpens.
  \item Training on forward data alone reduces $A$ and lets $B$ grow: the classic
        mode collapse.
  \item Since $\Delta = A + B - C$ contains both, a system in which one of the two
        data streams is inactive cannot reduce $\Delta$ past the term nothing is
        attacking, no matter how long it runs or how the remaining branch is
        weighted.
\end{itemize}
The last point is worth emphasising because it is easy to misdiagnose as slow
convergence rather than as a stalled component.

\subsection{What is estimable in practice}

\begin{table}[h]
\centering
\begin{tabular}{@{}lll@{}}
\toprule
Quantity & Estimator & Notes \\
\midrule
$\JF$ & mean of $\log w$ over forward rollouts & unbiased for $\log\Zstar - A$ \\
$\JB$ & mean of $\log w$ over buffer draws & may exceed $\log\Zstar$ \\
$\Delta$ & $\JB - \JF$ & $= A + B - C$ \\
$\log\Zstar$ & $\operatorname{logmeanexp}$ of $w$ on the \emph{forward} stream & Lemma~\ref{lem:is}; biased low \\
$A$ & $\operatorname{logmeanexp}(\log w) - \operatorname{mean}(\log w)$, forward & the ``Jensen gap'' \\
$C$ & --- & not identifiable without $\pi$ \\
\bottomrule
\end{tabular}
\end{table}

\noindent
Two cautions. First, the $\operatorname{logmeanexp}$ estimator of $\log\Zstar$ is
biased low by Jensen and has variance growing with $\operatorname{Var}(\log w)$,
precisely the regime where one wants it most; the resulting estimate of $A$ is
therefore an underestimate. Second, $B$ and $C$ are not separately identifiable
from $\JF$ and $\JB$ alone---only the combination $C - B$ is, since
$\log\Zstar - \JB = C - B$. Separating them requires an independent handle on the
buffer, e.g.\ knowing $\mu$ was drawn from a thermalised procedure.

\section{Summary}

\begin{align*}
  \JF &= \log\Zstar - A                 && A = \KL{\PF}{\pi\PB} && \text{(policy: mode-seeking)}\\
  \JB &= \log\Zstar - C + B             && B = \KL{\mu\PB}{\PF} && \text{(policy: mass-covering)}\\
  \Delta &= A + B - C                   && C = \KL{\mu}{\pi}    && \text{(buffer quality)}
\end{align*}

\noindent
Reading the identity as a design principle:
\begin{enumerate}
  \item \textbf{Objective:} $A$ down. It bounds the terminal-marginal error
        (Proposition~\ref{prop:dpi}) and nothing else on this list does.
  \item \textbf{Constraint:} $B$ bounded. Losing coverage is not recoverable by
        further optimisation of $A$.
  \item \textbf{Precondition:} $C$ small. No gradient reduces it; it is a
        data-preparation problem.
  \item \textbf{Nuisance:} $Z$. Maintain it well enough that neither residual is
        systematically corrupted, and never below $\mathrm{ELBO}(\mu)$
        (Corollary~\ref{cor:impossible}).
  \item \textbf{Diagnostic, never a target:} $\Delta$, which improves when the
        buffer degrades.
\end{enumerate}

\end{document}
